% !TeX root = ../papers/nc-p3.tex
\section{符号付き判別被覆上の対根と節点の小解消}
\label{sec:paired-root-resolution}
通常倍率を一組固定し、重み$(1,1,1,1,2)$の射影空間に
$X:\ S^2=F(a,b,c,d)$を取る。$D_0\ne0$では
$s=S/(\lambda(z^2-1)D_0)$として$s^2=\delta$である。
共通倍率を局所的に一つの非零成分で固定し、根$x,y$を各々射影直線に完備化する。
規格化を変えると根は共通倍率の二乗で変換されるので、これらの射影直線束は貼り合う。
\begin{theorem}\label{thm:paired-root-small-resolution}
$X$の特異点は十六の通常三次元二重点である。$D_0\ne0$の部分で、
符号$s,-s$に対応する二つの根のグラフ閉包は射影的小解消を与える。
八つの節点を除けばこれは同型で、各例外ファイバーは正規束
$\mathcal O_{\mathbb P^1}(-1)^{\oplus2}$を持つ一本の射影直線である。
その第二行列$[A_1]$は対応する接触二次曲線の全体を一対一にパラメータ付ける。
被覆の符号反転は各節点で二つの小解消を交換する。
\end{theorem}
\begin{proof}
重み付き射影空間の頂点は$X$に属さない。分岐面の滑らかな点では$X$も滑らかで、
面の節点の非退化二次項に$S^2$を加えると三次元通常二重点となる。
既約な分岐四次の二重被覆で、特異点は孤立するから$X$は整かつ正規である。

$[U:V]$を$x$の斉次根とし、
\[
 M(s)=\begin{pmatrix}2p&m-s\\m+s&2n\end{pmatrix},
 \qquad M(s)\binom UV=0
\]
を取る。$\det M=s^2-\delta$で、$V\ne0$なら一行目が$2px+m=s$、
二行を合わせて$px^2+mx+n=0$となる。非零行列では核直線は一意に正則に定まる。
無限遠根や最高次係数の零点でも同じである。零行列の点は$P\equiv0$であり、
通常倍率では定理\ref{thm:klein-all-pairings}と前の片側零の排除により八節点で尽きる。
そこで$p,m,n$の微分は独立である。実際、$\delta=m^2-4pn$の二次項が非退化で、
この三つの一次項が従属ならその階数は三に達しない。
従ってこの三つはエタール局所座標であり、$s$を加えた四座標から
$M$の四成分への変換は線型同型になる。

階数一以下の二行二列行列の核直線を指定した空間は、射影直線上で各行が
指定直線を消す共変ベクトルの対、すなわち$\mathcal O(-1)^{\oplus2}$の全空間である。
ゆえにこの射影的な核の表示は滑らかで、小解消となる。
$Q$について$-s$を使っても同様である。両根の被約グラフ閉包から第一の小解消への
射影は固有かつ双有理なので全射である。節点上では全ての点が既知の
$H=xy+Bx+Cy+D=0$の双斉次化を満たす。$BC-D$は八表示の全てで単元である。
実際、分子分母の非定数因子は
$z,\lambda,z\pm1,\lambda^2z-1,\lambda\pm1,\lambda z\pm1$
に限られる。よってその曲線は射影直線の自己同型のグラフであり、どの第一根にも
第二根は高々一つである。射影は準有限かつ固有なので有限、正規な第一の小解消への
有限双有理射なので同型となる。第二根への射影も同様である。

\path{computations/verify_paired_root_resolution.py}は各節点で二組のJacobian行列式と
$BC-D$が単元であることも検算する。節点の列を$N$とすると、
$H$から得た$A_1(x)$は$F(A_1)=0$と$(T^{-1}N)^t A_1=0$を満たす。
四成分の共通分母を払った二次以下の四式の$1,x,x^2$係数行列は、
通常倍率で単元となる三次小行列を持つ。よって無限遠にも基点がなく、
射影化した表示は次数二の完全線型系である。その像が上記の接触二次曲線である。
この完備化のうち元の可逆窓に戻る点は全て長さ6の障害を持つ。

最後に$M(-s)=M(s)^t$である。従って符号反転は核直線と左核直線を交換する。
固定した右核を持つ階数一行列をスカラー倍で零へ近づけると、
左核は任意に選べる。よって符号反転は同じ小解消の例外直線を越える正則な自己写像にはならず、
二つの小解消の間の通常のフロップを与える。
\end{proof}
